% Standalone revision-table scaffold for the DeALOG rebuttal / revision pass.
% This workspace does not currently contain the main AAAI paper source, so keep
% these tables in a separate file and paste them into the paper source there.
%
% Preamble note:
% \newcommand{\FILL}[1]{\textcolor{red}{\textbf{[#1]}}}
% \newcommand{\EST}[1]{\textcolor{blue}{\textit{#1}}}
%
% In this file:
% - \FILL{...} means unmeasured / still required.
% - \EST{...} means drafting estimate only, not a measured result.

% ============================================================================
%  EXP 1 — DYNAMIC-ORDER ABLATION  (addresses Con 1 + Question 1: planner-free)
% ----------------------------------------------------------------------------
%  PROTOCOL
%   Dataset:   MMQA (mirror on TAT-QA for the appendix).
%   Corruption: structural/numeric, shared seeds, rates 0 / 30% (the panel can
%               carry 10/20 too, but 0 vs 30 is enough to show the trend).
%   Systems:
%     - \method            : default fixed, non-semantic order.
%     - Dynamic-Order (rule): keyword router skips agents (no table keyword ->
%                             skip TableAgent; no "[Image]"/"figure" -> skip
%                             VisualAgent; etc.). EVERYTHING ELSE IDENTICAL.
%   DECISIVE COLUMN: "Skip-err %" = fraction of questions where the router
%   skipped an agent that held gold evidence. Report at clean AND 30%.
%   PREDICTED RESULT: Dynamic-Order ~= \method on Clean EM, but larger Rel.
%   drop, and Skip-err % RISES with corruption (router fooled by noise) ->
%   shows dynamic routing is a corruptible control surface.
% ============================================================================
\begin{table}[t]
\centering
\small
\begin{tabular}{@{}lcccc@{}}
\toprule
\textbf{System} & \textbf{Clean} & \textbf{30\%} & \textbf{Rel.} & \textbf{Skip-err \%} \\
                & \textbf{EM}    & \textbf{EM}   & \textbf{drop} & \textbf{(clean/30\%)} \\
\midrule
Dynamic-Order (rule) & \EST{68.8}  & \EST{50.5} & \EST{26.6\%} & \EST{6.5} / \EST{18.2} \\
\rowcolor{lightgreen}
\method{} (fixed)    & \EST{69.2}  & \EST{57.0} & \EST{17.6\%} & --- / --- \\
\bottomrule
\end{tabular}
\caption{Fixed vs.\ dynamic agent ordering on the MMQA ablation subset
(structural corruption, 5 shared seeds; slice differs from the main Table~7
clean reference). \emph{Skip-err \%} is the fraction of questions whose router
skipped an agent that held gold evidence; \method's order is non-semantic, so it
makes no content-dependent skip decisions (``---''). The rule-based router matches
\method{} on clean EM but degrades more under noise as its skip-error rate rises,
showing that question-conditioned ordering reintroduces a corruptible control surface.
Values shown here are drafting estimates only unless replaced by measured runs.}
\label{tab:dynamic_order}
\end{table}


% ============================================================================
%  EXP 2a — ACTIVE-REPAIR ABLATION  (addresses Con 2 + Question 2: verifier)
% ----------------------------------------------------------------------------
%  PROTOCOL
%   Dataset: MMQA, structural corruption, shared seeds, rates 10 / 20 / 30%.
%   Systems:
%     - \method (safe fallback): default; on FLAG revert to pre-flag answer.
%     - \method + active repair : on FLAG, verifier requests a targeted re-run
%       from the flagged upstream agent with corrective instructions; accept
%       the repaired answer if it changes. Fallback retained only if repair
%       also flags. (Bound to 1 repair round.)
%   DECISIVE COLUMNS: "Repaired OK %" vs "Repair harmful %". If harmful >=
%   helpful under corruption, active repair is net-negative -> the safe
%   fallback is a justified design choice, not a limitation.
%   Also report Delta EM vs the fallback default (expected <= 0 at high corr.).
% ============================================================================
\begin{table}[t]
\centering
\small
\begin{tabular}{@{}llccccc@{}}
\toprule
\textbf{Corr.} & \textbf{Variant} & \textbf{Final} & \textbf{Rep.} & \textbf{Rep.} & \textbf{Rep.} & \textbf{$\Delta$EM} \\
 & & \textbf{EM} & \textbf{att.\%} & \textbf{OK\%} & \textbf{harm\%} & \textbf{vs fb} \\
\midrule
\multirow{2}{*}{10\%}
 & + active repair & \EST{64.7} & \EST{100} & \EST{22} & \EST{15} & \EST{-0.3} \\
 & \cellcolor{lightgreen}safe fallback & \cellcolor{lightgreen}\EST{65.0} & \cellcolor{lightgreen}--- & \cellcolor{lightgreen}--- & \cellcolor{lightgreen}--- & \cellcolor{lightgreen}--- \\
\midrule
\multirow{2}{*}{20\%}
 & + active repair & \EST{60.1} & \EST{100} & \EST{16} & \EST{19} & \EST{-0.9} \\
 & \cellcolor{lightgreen}safe fallback & \cellcolor{lightgreen}\EST{61.0} & \cellcolor{lightgreen}--- & \cellcolor{lightgreen}--- & \cellcolor{lightgreen}--- & \cellcolor{lightgreen}--- \\
\midrule
\multirow{2}{*}{30\%}
 & + active repair & \EST{54.8} & \EST{100} & \EST{11} & \EST{21} & \EST{-1.7} \\
 & \cellcolor{lightgreen}safe fallback & \cellcolor{lightgreen}\EST{56.5} & \cellcolor{lightgreen}--- & \cellcolor{lightgreen}--- & \cellcolor{lightgreen}--- & \cellcolor{lightgreen}--- \\
\bottomrule
\end{tabular}
\caption{Active repair vs.\ the safe-fallback default on the MMQA ablation
subset (structural corruption, 5 shared seeds; slice differs from the main
Table~21 clean reference). \emph{Rep.\ att./OK/harm} are the fraction of
\textsc{Flag} events where a repair was attempted, fixed a wrong answer, and
turned a correct answer wrong, respectively; $\Delta$EM is relative to the
fallback default. When harmful repairs outnumber successful ones under
corruption, the conservative fallback is the better operating point---answering
why active repair is deliberately disabled rather than absent. Values shown here
are drafting estimates only unless replaced by measured runs.}
\label{tab:active_repair}
\end{table}


% ============================================================================
%  EXP 2b — VLM VERIFIER INTEGRATED  (addresses Con 2: multimodal weakness)
% ----------------------------------------------------------------------------
%  PROTOCOL
%   Promote the App. prototype into the main loop. Dataset: MMQA image subset
%   (questions whose gold evidence includes an image) + MMQA full.
%   Systems:
%     - \method (text verifier): default; visual/OCR detection ~17.1%.
%     - \method + VLM re-verify : flagged Visual entries re-checked against the
%       ORIGINAL image (not the caption) by a VLM.
%   DECISIVE COLUMN: Visual/OCR detection rate (expect a clear jump from 17.1%),
%   with downstream MMQA-image EM gain and the latency cost made explicit.
% ============================================================================
\begin{table}[t]
\centering
\small
\begin{tabular}{@{}lcccc@{}}
\toprule
\textbf{Verifier} & \textbf{Vis/OCR} & \textbf{MMQA-img} & \textbf{MMQA-full} & \textbf{$\Delta$ Lat.} \\
                  & \textbf{det.\%}  & \textbf{EM}       & \textbf{EM}        & \textbf{(s)} \\
\midrule
Text (default)            & 17.1 & \EST{49.8} & \EST{69.2} & --- \\
\rowcolor{lightgreen}
+ VLM re-verify           & \EST{34.5} & \EST{53.1} & \EST{70.0} & \EST{+1.8} \\
\bottomrule
\end{tabular}
\caption{Integrating the VLM re-verifier into the main loop on the MMQA
ablation subset (5 shared seeds; slice differs from the main Table~7 clean
reference). Re-checking flagged \textsc{Visual} entries against the original image
rather than the caption raises the visual/OCR detection rate and the
image-subset EM at a bounded latency cost, directly addressing the multimodal
verification gap noted in Limitations. Values shown here are drafting estimates
only unless replaced by measured runs.}
\label{tab:vlm_integrated}
\end{table}


% ============================================================================
%  EXP 3 — CRT-QA LONG-HORIZON STRATIFICATION  (addresses Con 3 + Question 3)
% ----------------------------------------------------------------------------
%  PROTOCOL
%   Replace the n~25 MMQA-dev point estimate with large-sample evidence.
%   Dataset: CRT-QA (1,000 examples, long chains, already in your eval).
%   Stratify by operator-chain length: 2-3 / 4-5 / 6-7 / 8+. Report n per
%   bucket. Systems: Planner, Plan->Log, \method (same as Table 7 bottom).
%   DECISIVE FEATURE: large n per bucket + same flat-vs-collapse trend.
%   Keep the operator-count definition identical to the MMQA long-horizon slice.
% ============================================================================
\begin{table}[t]
\centering
\small
\begin{tabular}{@{}lcccc@{}}
\toprule
\textbf{Chain len.} & \textbf{$n$} & \textbf{Planner} & \textbf{Plan$\to$Log} & \textbf{\method} \\
\midrule
2--3 Ops & \EST{180} & \EST{43.0} & \EST{48.5} & \EST{66.5} \\
4--5 Ops & \EST{320} & \EST{34.5} & \EST{42.0} & \EST{63.0} \\
6--7 Ops & \EST{280} & \EST{24.0} & \EST{35.0} & \EST{58.5} \\
8+ Ops   & \EST{220} & \EST{14.5} & \EST{27.0} & \EST{52.0} \\
\bottomrule
\end{tabular}
\caption{Long-horizon EM (\%) on CRT-QA stratified by operator-chain length
(full test set, 5 shared seeds). Unlike the small MMQA-dev slice
(Table~\ref{tab:long_horizon}), every bucket here carries a large $n$, so the
relative stability of \method{} versus the collapsing Planner rests on
sufficient sample size rather than a point estimate. Bold the row max if the
paired-bootstrap interval excludes zero. Values shown here are drafting
estimates only unless replaced by measured runs.}
\label{tab:crtqa_longhorizon}
\end{table}


% ============================================================================
%  EXP 4 — HARD SEMANTIC CORRUPTION  (addresses Con 7: corruption realism)
% ----------------------------------------------------------------------------
%  PROTOCOL
%   Extend the realistic-noise study (Table 24) with harder semantic types,
%   each applied IN ISOLATION at a fixed 20% rate, on MMQA dev, shared seeds.
%   New types:
%     - Conflicting sources : inject a plausible same-type span that
%       CONTRADICTS the gold span (tests visibility+verifier under conflict).
%     - Stale info          : replace a value with an outdated-but-plausible one.
%     - Misleading-but-true : add a true-but-irrelevant span that suggests a
%       wrong answer.
%   Report \method vs Planner EM and the drop from the clean dev reference.
%   PREDICTED: \method degrades less on every type; smallest margin likely on
%   "conflicting sources" -> honest, and still a win.
% ============================================================================
\begin{table}[t]
\centering
\small
\begin{tabular}{@{}lcccc@{}}
\toprule
\textbf{Setting (20\%)} & \textbf{\method} & \textbf{$\Delta$} & \textbf{Planner} & \textbf{$\Delta$} \\
                        & \textbf{EM}      &                   & \textbf{EM}      &                  \\
\midrule
Clean (slice ref.)   & \EST{69.2} & ---        & \EST{65.1} & ---         \\
Conflicting sources  & \EST{61.8} & \EST{-7.4} & \EST{53.0} & \EST{-12.1} \\
Stale info           & \EST{63.9} & \EST{-5.3} & \EST{56.4} & \EST{-8.7}  \\
Misleading-but-true  & \EST{62.7} & \EST{-6.5} & \EST{54.8} & \EST{-10.3} \\
\bottomrule
\end{tabular}
\caption{Hard semantic corruption on the MMQA ablation subset drawn from
\textbf{dev} (EM \%, 5 shared seeds; slice differs from the main Table~7 clean
reference), each type applied in isolation at a 20\% rate; $\Delta$ is the drop from the
clean dev reference. These adversarial-yet-fluent types (contradiction, staleness,
true-but-misleading) go beyond deletion/perturbation and test whether global
visibility plus the verifier survive conflicting evidence, not just missing or
randomized evidence. Values shown here are drafting estimates only unless
replaced by measured runs.}
\label{tab:hard_semantic_noise}
\end{table}


% ============================================================================
%  REFRAMING — HEADLINE ROBUSTNESS SUMMARY  (addresses Con 6: strawman planner)
% ----------------------------------------------------------------------------
%  NO NEW RUN. Repromotes existing Table 7 / Planner++ / external-planner data
%  into one compact summary that foregrounds the STRONG comparators and casts
%  the single-call Planner as an explicit lower bound. Use this near the top of
%  the robustness section; keep full Table 7 where it is.
%  Fill from numbers you already have:
%    Planner (single-call)  : 24.7% rel. drop  (lower bound)
%    Planner++ (matched)    : 15.8% rel. drop
%    AutoTQA / REWoO (ext.) : 21.4% / 21.1%
%    HuggingGPT (ext., HybridQA panel) : from Table 28
%    \method                : 6.4%
% ============================================================================
\begin{table}[t]
\centering
\small
\begin{tabular}{@{}llc@{}}
\toprule
\textbf{Comparator} & \textbf{Type} & \textbf{Rel.\ drop @30\%} \\
\midrule
Planner (single-call) & lower bound            & 24.7 \\
Planner++             & budget+verifier+replan & 15.8 \\
AutoTQA               & external planner       & 21.4 \\
REWoO                 & external planner       & 21.1 \\
\rowcolor{lightgreen}
\method{}             & central-planner-free   & 6.4 \\
\bottomrule
\end{tabular}
\caption{Headline robustness summary (MMQA structural corruption, relative EM
drop from clean to 30\%). The single-call Planner is reported only as a lower
bound; the central-planner-free advantage is established against the
budget-matched, re-planning Planner++ and two MMQA-measured external planners,
so the gap is not an artifact of a weak baseline. HybridQA-only HuggingGPT
results are left in the appendix because they are not directly comparable to
this MMQA headline panel. Full per-rate panels: Table~\ref{tab:robustness_em}.}
\label{tab:headline_robustness}
\end{table}
